Autoregressive transformer models often generate false information with high confidence, a failure mode known as hallucination. We introduce the \textit{Safety Ratio} (SR), a simple mechanistic metric that quantifies the balance between a model's grounding in local statistics (Layer 0 ``Sentinel'' neurons) and its semantic completion processing (Layers 2--4 ``Engine'' neurons). Using 301 prompts across three categories (truth, nonsense, adversarial) on Qwen2.5-1.5B, we show that a high-recall SR threshold (SR $< 0.368$, chosen on a validation split as the 95th percentile of adversarial SR) detects 95\% of adversarial inputs on a held-out test set, with statistically significant separation from truthful prompts ($p = 0.0052$, Cohen's $d = 0.40$) and improved recall over perplexity-based baselines (95\% vs 44\%). Adversarial prompts exhibit a distinct mechanistic signature---simultaneously lowest mean SR and lowest variance---consistent with robust Sentinel--Engine decoupling. We further conduct calibrated simulated scaling experiments on GPT-2-XL, Pythia-1.4B, and a 7B model, together with preliminary real runs on GPT-2 Small and GPT-2 Medium, to explore how this signature might behave across architectures; these analyses are hypothesis-generating and do not constitute definitive cross-model validation. Overall, SR provides a lightweight, interpretable signal for high-recall adversarial pre-filtering in cascade detection systems, trading precision for sensitivity in a way that is appropriate for safety-critical deployments.

